% GENERATED FILE. Edit canonical lessons, then run npm run generate:books.
% canonical-source-hash: fnv1a64:ef85cf3ee178a97a
% canonical-lessons: GU-C544-pher, GU-C544-taphavat, GU-C544-alag, GU-C544-nakal, GU-C544-barabar

\chapter{Same or Different}
\label{ch:gu-same-or-different}
\hlchaptermodality{544}

\begin{chapteropening}
\textbf{By the end of this chapter:} \emph{I can compare two things and say how they differ.}

Complete the last lesson of chapter 544: I can compare two things and say how they differ.
\end{chapteropening}

\section[\texorpdfstring{\gu{ફેર} (pher)}{ફેર (pher)}]{\gu{ફેર} (pher) — a difference}

\label{lesson:GU-C544-pher}



\begin{quote}
Before the new one: say the Gujarati for capable, then the Gujarati for special.
\end{quote}

\subsection*{You'll want to know: \gu{ફેર}}
\textbf{\gu{ફેર}} — \emph{pher} — \textquotedblleft{}a difference\textquotedblright{}.

\begin{quote}
\gu{જો} \gu{તમે} \gu{બંને} \gu{ગીતો} \gu{સાંભળશો} \gu{તો} \gu{તમે} \gu{ફેર} \gu{ઓળખી} \gu{શકશો}.
\end{quote}

\begin{quote}
\emph{jo tame banne gīto sā̃bhaḷśo to tame pher oḷkhī śakśo.}
\end{quote}

\begin{quote}
If you listen to both songs, you'll be able to tell the difference.
\end{quote}

A masculine noun.

\begin{cousinweb}
\textbf{\gu{ફેર}} (\emph{pher}) is filed by CDIAL under \emph{*phirati}, \textquotedblleft{}turns\textquotedblright{}, with a Gujarati \emph{pher}, \textquotedblleft{}a turn\textquotedblright{} (CDIAL 9078); how the word came to mean a difference is not explained there.
\end{cousinweb}

\begin{grammarlens}[title={What you've built}]
One more word for saying what happens if.
\end{grammarlens}

\subsection*{Your turn}
\begin{itemize}
\raggedright
  \item \emph{Say it:} \emph{pher}
  \item \emph{Say it:} \emph{pher}, once more
  \item \emph{Say it:} \emph{jo tame banne gīto sā̃bhaḷśo to tame pher oḷkhī śakśo.} — If you listen to both songs, you'll be able to tell the difference.
  \item \emph{Recall:} say the Gujarati for capable, then the Gujarati for special, then say \emph{pher} again
\end{itemize}

\begin{tcolorbox}[breakable,colback=teal!4,colframe=teal!35!black,title={Before you move on}]
What does \gu{ફેર} mean? (\textquotedblleft{}A difference\textquotedblright{}.) Say it once more.
\end{tcolorbox}
\section[\texorpdfstring{\gu{તફાવત} (taphāvat)}{તફાવત (taphāvat)}]{\gu{તફાવત} (taphāvat) — a distinction}

\label{lesson:GU-C544-taphavat}



\begin{quote}
Before the new one: say the Gujarati for special, then the Gujarati for a difference.
\end{quote}

\subsection*{You'll want to know: \gu{તફાવત}}
\textbf{\gu{તફાવત}} — \emph{taphāvat} — \textquotedblleft{}a distinction\textquotedblright{}.

\begin{quote}
\gu{જો} \gu{બંનેના} \gu{ભાવમાં} \gu{મોટો} \gu{તફાવત} \gu{હશે} \gu{તો} \gu{હું} \gu{સસ્તું} \gu{ખરીદીશ}.
\end{quote}

\begin{quote}
\emph{jo bannenā bhāvmā̃ moṭo taphāvat haśe to hũ sastũ kharīdīś.}
\end{quote}

\begin{quote}
If there's a big difference between the two prices, I'll buy the cheaper one.
\end{quote}

A masculine noun.

\begin{cousinweb}
\textbf{\gu{તફાવત}} (\emph{taphāvat}) is usually traced to Arabic \emph{tafāwut}, \textquotedblleft{}a difference, a disparity\textquotedblright{}, through Persian.
\end{cousinweb}

\begin{grammarlens}[title={What you've built}]
One more word for saying what happens if.
\end{grammarlens}

\subsection*{Your turn}
\begin{itemize}
\raggedright
  \item \emph{Say it:} \emph{taphāvat}
  \item \emph{Say it:} \emph{taphāvat}, once more
  \item \emph{Say it:} \emph{jo bannenā bhāvmā̃ moṭo taphāvat haśe to hũ sastũ kharīdīś.} — If there's a big difference between the two prices, I'll buy the cheaper one.
  \item \emph{Recall:} say the Gujarati for special, then the Gujarati for a difference, then say \emph{taphāvat} again
\end{itemize}

\begin{tcolorbox}[breakable,colback=teal!4,colframe=teal!35!black,title={Before you move on}]
What does \gu{તફાવત} mean? (\textquotedblleft{}A distinction\textquotedblright{}.) Say it once more.
\end{tcolorbox}
\section[\texorpdfstring{\gu{અલગ} (alag)}{અલગ (alag)}]{\gu{અલગ} (alag) — separate}

\label{lesson:GU-C544-alag}



\begin{quote}
Before the new one: say the Gujarati for a difference, then the Gujarati for a distinction.
\end{quote}

\subsection*{You'll want to know: \gu{અલગ}}
\textbf{\gu{અલગ}} — \emph{alag} — \textquotedblleft{}separate\textquotedblright{}.

\begin{quote}
\gu{જો} \gu{બાળકો} \gu{અલગ} \gu{ઓરડામાં} \gu{સૂઈ} \gu{જશે} \gu{તો} \gu{ઘર} \gu{શાંત} \gu{રહેશે}.
\end{quote}

\begin{quote}
\emph{jo bāḷko alag orḍāmā̃ sūī jaśe to ghar śānt raheśe.}
\end{quote}

\begin{quote}
If the children sleep in a separate room, the house will stay quiet.
\end{quote}

\begin{cousinweb}
\textbf{\gu{અલગ}} (\emph{alag}) is usually linked to Sanskrit \emph{alagna}, \textquotedblleft{}unconnected\textquotedblright{} (CDIAL 700), under which CDIAL gives a Gujarati \emph{aḷgũ}, \textquotedblleft{}not touching\textquotedblright{}; this form is not listed there.
\end{cousinweb}

\begin{grammarlens}[title={What you've built}]
One more word for saying what happens if.
\end{grammarlens}

\subsection*{Your turn}
\begin{itemize}
\raggedright
  \item \emph{Say it:} \emph{alag}
  \item \emph{Say it:} \emph{alag}, once more
  \item \emph{Say it:} \emph{jo bāḷko alag orḍāmā̃ sūī jaśe to ghar śānt raheśe.} — If the children sleep in a separate room, the house will stay quiet.
  \item \emph{Recall:} say the Gujarati for a difference, then the Gujarati for a distinction, then say \emph{alag} again
\end{itemize}

\begin{tcolorbox}[breakable,colback=teal!4,colframe=teal!35!black,title={Before you move on}]
What does \gu{અલગ} mean? (\textquotedblleft{}Separate\textquotedblright{}.) Say it once more.
\end{tcolorbox}
\section[\texorpdfstring{\gu{નકલ} (nakal)}{નકલ (nakal)}]{\gu{નકલ} (nakal) — a copy (an imitation)}

\label{lesson:GU-C544-nakal}



\begin{quote}
Before the new one: say the Gujarati for a distinction, then the Gujarati for separate.
\end{quote}

\subsection*{You'll want to know: \gu{નકલ}}
\textbf{\gu{નકલ}} — \emph{nakal} — \textquotedblleft{}a copy (an imitation)\textquotedblright{}.

\begin{quote}
\gu{જો} \gu{આ} \gu{ચિત્ર} \gu{નકલ} \gu{હશે} \gu{તો} \gu{તેની} \gu{કિંમત} \gu{ઓછી} \gu{હશે}.
\end{quote}

\begin{quote}
\emph{jo ā citra nakal haśe to tenī kimmat ochī haśe.}
\end{quote}

\begin{quote}
If this painting is a copy, it will be worth less.
\end{quote}

A feminine noun.

\begin{cousinweb}
\textbf{\gu{નકલ}} (\emph{nakal}) is usually traced to Arabic \emph{naql}, \textquotedblleft{}a transcript, a copy\textquotedblright{}, which reached Gujarati through Persian.
\end{cousinweb}

\begin{grammarlens}[title={What you've built}]
One more word for saying what happens if.
\end{grammarlens}

\subsection*{Your turn}
\begin{itemize}
\raggedright
  \item \emph{Say it:} \emph{nakal}
  \item \emph{Say it:} \emph{nakal}, once more
  \item \emph{Say it:} \emph{jo ā citra nakal haśe to tenī kimmat ochī haśe.} — If this painting is a copy, it will be worth less.
  \item \emph{Recall:} say the Gujarati for a distinction, then the Gujarati for separate, then say \emph{nakal} again
\end{itemize}

\begin{tcolorbox}[breakable,colback=teal!4,colframe=teal!35!black,title={Before you move on}]
What does \gu{નકલ} mean? (\textquotedblleft{}A copy (an imitation)\textquotedblright{}.) Say it once more.
\end{tcolorbox}
\section[\texorpdfstring{\gu{બરાબર} (barābar)}{બરાબર (barābar)}]{\gu{બરાબર} (barābar) — evenly matched}

\label{lesson:GU-C544-barabar}



\begin{quote}
Before the new one: say the Gujarati for separate, then the Gujarati for a copy.
\end{quote}

\subsection*{You'll want to know: \gu{બરાબર}}
\textbf{\gu{બરાબર}} — \emph{barābar} — \textquotedblleft{}evenly matched\textquotedblright{}.

\begin{quote}
\gu{જો} \gu{બંને} \gu{ભાગ} \gu{બરાબર} \gu{હશે} \gu{તો} \gu{કોઈ} \gu{ફરિયાદ} \gu{નહીં} \gu{કરે}.
\end{quote}

\begin{quote}
\emph{jo banne bhāg barābar haśe to koī phariyād nahī̃ kare.}
\end{quote}

\begin{quote}
If both parts are equal, nobody will complain.
\end{quote}

\begin{cousinweb}
\textbf{\gu{બરાબર}} (\emph{barābar}) is usually traced to Persian \emph{barābar}, \textquotedblleft{}equal, level\textquotedblright{}. It is also an everyday \textquotedblleft{}right, exactly\textquotedblright{}: \textbf{\gu{બરાબર} \gu{છે}}.
\end{cousinweb}

\begin{grammarlens}[title={What you've built}]
One more word for saying what happens if.
\end{grammarlens}

\subsection*{Your turn}
\begin{itemize}
\raggedright
  \item \emph{Say it:} \emph{barābar}
  \item \emph{Say it:} \emph{barābar}, once more
  \item \emph{Say it:} \emph{jo banne bhāg barābar haśe to koī phariyād nahī̃ kare.} — If both parts are equal, nobody will complain.
  \item \emph{Recall:} say the Gujarati for separate, then the Gujarati for a copy, then say \emph{barābar} again
\end{itemize}

\begin{tcolorbox}[breakable,colback=teal!4,colframe=teal!35!black,title={Before you move on}]
What does \gu{બરાબર} mean? (\textquotedblleft{}Evenly matched\textquotedblright{}.) Say it once more.
\end{tcolorbox}
